\section{Reactions}
Sometimes called \emph{Reaction mechanism}.

If the experiment involves chemical reactions, write them down and explain them: this is one of the key parts of a chemistry report. Use the \verb|\ce{...}| command of the \texttt{mhchem} package for formulae and equations, and the \texttt{equation} environment to number them.

\begin{equation}
    \ce{2 H2O2 -> 2 H2O + O2 ^}
    \label{eq:peroxide}
\end{equation}

Some \texttt{mhchem} examples:
\begin{itemize}
    \item ions and charges: \ce{SO4^2-}, \ce{[Cu(NH3)4]^2+};
    \item equilibrium: \ce{H2O <=> H+ + OH-};
    \item bonds: \ce{HC#CH}, \ce{H2C=CH2}, \ce{H3C-CH3};
    \item conditions above the arrow: \ce{A ->[$\Delta$] B};
    \item states of matter and precipitates: \ce{Ag+(aq) + Cl-(aq) -> AgCl(s) v}.
\end{itemize}

Refer to numbered reactions with \verb|\cref|: the decomposition in \cref{eq:peroxide} is catalysed by \ce{MnO2}.
